\chapter{Fungsi Baku}\label{ch:b1:functions}

Analisis hanya seberguna persediaan fungsi yang kita kuasai. Pada
koleksi warisan sekolah menengah --- pangkat, eksponensial,
logaritma, fungsi trigonometri --- bab ini menambahkan fungsi
inversnya ($\arcsin$, $\arccos$, $\arctan$) dan keluarga hiperbolik.
Turunan dipakai dengan bebas pada taraf sekolah menengah; teori
di balik fungsi invers dituntaskan pada
\cref{ch:b1:continuity,ch:b1:derivative}.

\section{Eksponensial, logaritma, pangkat}

\begin{proposition}[Ringkasan dan pencirian]\label{prop:b1:functions:expln}
$\exp \colon \R \to \intoo{0}{+\infty}$ dan $\ln \colon
\intoo{0}{+\infty} \to \R$ adalah dua bijeksi yang saling invers, naik
tegas, dengan
\[
\exp(x + y) = \exp x \exp y, \qquad \ln(xy) = \ln x + \ln y,
\qquad \exp' = \exp, \qquad \ln'(x) = \frac 1x .
\]
Lebih lanjut, $\exp$ adalah \emph{satu-satunya} fungsi $f \colon
\R \to \R$ yang dapat diturunkan dengan $f' = f$ dan $f(0) = 1$.
\end{proposition}

\begin{proof}
Ringkasannya adalah bahan sekolah menengah. Untuk ketunggalannya, misalkan $f' = f$,
$f(0) = 1$, lalu tulis $g(x) = f(x)\,\eu^{-x}$. Maka $g' = f'\eu^{-x} -
f\eu^{-x} = 0$, jadi $g$ konstan dan sama dengan $g(0) = 1$: $f = \exp$.
\end{proof}

\begin{definition}[Pangkat umum]\label{def:b1:functions:powers}
Untuk $x > 0$ dan $\alpha \in \R$:
$\;x^\alpha = \eu^{\alpha \ln x}$\index{fungsi pangkat}. Untuk $a > 0$,
$a \neq 1$, \emph{logaritma dengan bilangan pokok $a$} adalah $\log_a x =
\frac{\ln x}{\ln a}$, yaitu invers $x \mapsto a^x$.
\end{definition}

\begin{example}[Menyelesaikan persamaan eksponensial]\label{ex:b1:functions:expequations}
Selesaikan $2^x = 5^{\,x-1}$ di $\R$. Kedua ruasnya positif, jadi ambil
logaritmanya --- langkah yang dapat dibalik:
\[
x\ln2 = (x - 1)\ln5
\iff x(\ln2 - \ln5) = -\ln5
\iff x = \frac{\ln5}{\ln5 - \ln2} = \frac{\ln 5}{\ln\frac52}
\approx 1.756 .
\]
Berikutnya, selesaikan $x^{\sqrt2} = 3$ untuk $x > 0$: pangkatkan dengan
$\frac1{\sqrt2}$ (yakni terapkan bijeksi inversnya):
$x = 3^{1/\sqrt2} = \eu^{(\ln 3)/\sqrt2} \approx 2.175$. Inti
gagasannya: setiap persamaan yang mencampur pangkat terurai lewat $\ln$ dan
$\exp$, karena definisi $x^\alpha = \eu^{\alpha\ln x}$
menyusutkan semua manipulasi pangkat menjadi aritmetika eksponen ---
tetapi hanya pada daerah $x > 0$ yang di situlah definisi itu hidup.
\end{example}

\begin{example}[Waktu penggandaan]\label{ex:b1:functions:doubling}
Sebuah besaran tumbuh $3\%$ per langkah: setelah $n$ langkah ia
terkali $(1.03)^n$. Kapan ia berlipat dua? Selesaikan $(1.03)^n
\geq 2$:
\[
n\ln(1.03) \geq \ln2
\iff n \geq \frac{\ln 2}{\ln 1.03}
= \frac{0.6931}{0.02956} \approx 23.4 ,
\]
jadi penggandaan pertama terjadi pada langkah $24$. (``Aturan $72$''
para pemodal, yang menaksir waktu penggandaan sebagai
$72$ dibagi lajunya dalam persen, adalah perhitungan ini dengan
hampiran $\ln(1 + x) \approx x$, yang dikuantifikasi pada
\cref{ch:b1:taylor}.) Proses eksponensial
paling baik dinalar lewat logaritmanya: pada
skala itu, pertumbuhannya linear dan pertanyaannya menjadi pembagian.
\end{example}

\begin{example}[Sepanjang apa $2^{2026}$?]\label{ex:b1:functions:digits}
Banyaknya angka desimal sebuah bilangan bulat $N \geq 1$ adalah
$\floor{\log_{10} N} + 1$ (memang $N$ mempunyai $d$ angka tepat ketika
$10^{d-1} \leq N < 10^d$, yakni $d - 1 \leq \log_{10}N < d$). Untuk
$N = 2^{2026}$:
\[
\log_{10} 2^{2026} = 2026\,\log_{10}2
= 2026 \times 0.301030 = 609.887 ,
\]
jadi $2^{2026}$ mempunyai $610$ angka. Bagian pecahannya membawa
bonus: $10^{0.887} \approx 7.7$, jadi bilangan itu \emph{diawali}
$7$. Satu perkalian menjawab pertanyaan tentang bilangan yang tak
seorang pun akan pernah menuliskannya --- logaritma memampatkan ukuran
perkalian menjadi ukuran penjumlahan, dan itulah seluruh maknanya secara historis
(\cref{ex:b1:structures:expmorphism} membuat kalimat itu menjadi tepat).
\end{example}

\begin{proposition}[Kaidah pangkat dan perbandingan pertumbuhan]\label{prop:b1:functions:powerrules}
Untuk $x, y > 0$ dan $\alpha, \beta \in \R$:
\[
x^{\alpha+\beta} = x^\alpha x^\beta, \quad
(x^\alpha)^\beta = x^{\alpha\beta}, \quad
(xy)^\alpha = x^\alpha y^\alpha, \quad
(x^\alpha)' = \alpha\, x^{\alpha - 1}.
\]
Skala pertumbuhan\index{perbandingan pertumbuhan} ketika $x \to +\infty$, untuk setiap
$\alpha > 0$ dan $\beta > 0$:
\[
\frac{(\ln x)^\beta}{x^\alpha} \longrightarrow 0,
\qquad
\frac{x^\alpha}{\eu^{\beta x}} \longrightarrow 0 .
\]
\end{proposition}

\begin{proof}
Kesamaan itu menyalin kesamaan $\exp$ dan $\ln$ lewat
definisinya. Secara rinci untuk yang kedua (yang paling tak kasatmata): $x^\alpha
> 0$ dan $\ln(x^\alpha) = \alpha\ln x$ (terapkan $\ln$ pada
definisinya), sehingga
\[
(x^\alpha)^\beta = \eu^{\beta\ln(x^\alpha)}
= \eu^{\beta\alpha\ln x} = x^{\alpha\beta} ;
\]
yang pertama dan ketiga adalah penguraian dua baris yang sama lewat
$\exp(u + v) = \exp u\exp v$ dan $\ln(xy) = \ln x + \ln y$. Turunannya
adalah aturan rantai: $(\eu^{\alpha\ln x})' =
\frac{\alpha}{x} \eu^{\alpha\ln x}$.

Perbandingannya: dari $\frac{\ln t}{t} \to 0$ (yang dibuktikan pada jilid
sekolah menengah), substitusikan $t = x^{\alpha/\beta}$:
$\frac{\ln x}{x^{\alpha/\beta}} \to 0$, lalu pangkatkan dengan
$\beta$. Untuk yang kedua, substitusikan $x = \ln u$ pada yang pertama.
\end{proof}

\begin{example}[Dua limit yang harus dikuasai setiap pembaca]\label{ex:b1:functions:xx}
Berapa $\lim_{x \to 0^+} x^x$ dan $\lim_{x \to +\infty}
x^{1/x}$? Kedua pangkat itu terdefinisi lewat eksponensial, jadi
selesaikan dulu eksponennya:
\[
x^x = \eu^{x\ln x}, \qquad
x\ln x = -\frac{\ln(1/x)}{1/x} \longrightarrow 0
\quad (x \to 0^+),
\qquad\text{sehingga } x^x \longrightarrow \eu^0 = 1 ;
\]
dan $x^{1/x} = \eu^{(\ln x)/x} \to \eu^0 = 1$ ketika $x \to +\infty$,
langsung menurut \omterm{prop:b1:functions:powerrules}{perbandingan pertumbuhan}. Inti gagasannya: pangkat
tak tentu (yang berbentuk $0^0$ atau $\infty^0$) \emph{selalu} ditangani dengan
menuliskannya ulang $u^v = \eu^{v\ln u}$ lalu menganalisis hasil kali $v \ln u$
--- jangan pernah dengan menerka dari bilangan pokok dan eksponennya secara terpisah.
Fungsi $x \mapsto \frac{\ln x}{x}$ yang memutuskan kedua limit itu
dipelajari tuntas pada soal akhir pekan bab ini.
\end{example}

\begin{omfigure}
\begin{tikzpicture}
\begin{axis}[omaxis, width=9.5cm, height=6cm,
    xmin=0, xmax=5.2, ymin=-1.8, ymax=8.5,
    xtick={1,2,3,4,5}, ytick={2,4,6,8}]
  \addplot[omMeth, thick, domain=0.02:5, samples=200] {ln(x)}
    node[pos=0.97, below, font=\small] {$\ln x$};
  \addplot[omDef, thick, domain=0:5, samples=100] {sqrt(x)}
    node[pos=0.9, above, font=\small] {$\sqrt x$};
  \addplot[omProp, thick, domain=0:5] {x}
    node[pos=0.72, left, font=\small] {$x$};
  \addplot[omThm, thick, domain=0:2.14, samples=120] {exp(x)}
    node[pos=0.85, left, font=\small] {$\eu^x$};
\end{axis}
\end{tikzpicture}

{\small Skala pertumbuhan pada
\cref{prop:b1:functions:powerrules}: di dekat tepi kanan, $\ln x$
baru saja melewati $1.6$ sedangkan $\eu^x$ sudah meninggalkan bingkainya. Setiap
rasionya (logaritma dibagi pangkat, dan pangkat dibagi eksponensial) menuju $0$ ---
gambarnya hanya mengisyaratkan apa yang dibuat tepat oleh substitusi pada
buktinya.}
\end{omfigure}

\begin{remark}[Jebakan yang lazim dengan pangkat dan logaritma]\label{rem:b1:functions:pitfalls}
\begin{enumerate}
  \item \emph{Daerahnya.} $x^\alpha$ untuk $\alpha$ yang irasional menuntut
        $x > 0$; $(-8)^{1/3}$ sebaiknya dihindari dan diganti dengan ``akar
        pangkat tiga real dari $-8$'', karena kaidah
        $(x^\alpha)^\beta = x^{\alpha\beta}$ gugur diam-diam pada
        bilangan negatif: $\bigl((-8)^2\bigr)^{1/6} = 2 \neq -2$.
  \item \emph{$\sqrt{x^2} = \abs x$}, bukan $x$: melupakan
        nilai mutlaknya adalah sumber klasik hilangnya penyelesaian
        yang negatif.
  \item \emph{$\ln(xy) = \ln x + \ln y$ menuntut $x, y > 0$.} Untuk
        argumen yang negatif, $\ln(xy)$ bisa saja terdefinisi sedangkan
        ruas kanannya tidak.
  \item \emph{Fungsi mana yang tumbuh?} Pada $x^\alpha$ yang
        berubah adalah \emph{bilangan pokoknya}; pada $a^x$ yang berubah
        adalah \emph{eksponennya}. Pertumbuhan keduanya jauh berbeda
        (\cref{prop:b1:functions:powerrules}), dan ungkapan
        campuran seperti $x^{\ln x}$ atau $x^{1/x}$ harus
        ditulis ulang sebagai $\eu^{(\cdot)}$ sebelum dinalar ---
        seperti pada \cref{ex:b1:functions:xx}.
\end{enumerate}
\end{remark}

\section{Fungsi invers trigonometri}

\begin{definition}[$\arcsin$, $\arccos$, $\arctan$]\label{def:b1:functions:arc}
Pembatasan
\[
\sin \colon \intcc{-\tfrac\pi2}{\tfrac\pi2} \to \intcc{-1}{1},
\qquad
\cos \colon \intcc{0}{\pi} \to \intcc{-1}{1},
\qquad
\tan \colon \intoo{-\tfrac\pi2}{\tfrac\pi2} \to \R
\]
adalah bijeksi yang monoton tegas. \omterm{def:b1:logic:map}{Pemetaan} inversnya ditulis
$\arcsin$\index{arkus sinus}, $\arccos$\index{arkus kosinus} dan
$\arctan$\index{arkus tangen}. Jadi, misalnya, $y = \arcsin x$
($x \in \intcc{-1}{1}$) adalah \emph{satu-satunya} sudut di
$\intcc{-\frac\pi2}{\frac\pi2}$ yang sinusnya $x$.
\end{definition}

\begin{omfigure}
\begin{tikzpicture}
\begin{axis}[omaxis, width=5.6cm, height=5.6cm,
    xmin=-1.25, xmax=1.25, ymin=-1.8, ymax=1.8,
    xtick={-1,1}, ytick={-1.5708,1.5708},
    yticklabels={$-\frac{\pi}{2}$,$\frac{\pi}{2}$}]
  \addplot[omDef, very thick, domain=-1:1, samples=200]
    {rad(asin(x))};
\end{axis}
\begin{scope}[xshift=5.3cm]
\begin{axis}[omaxis, width=5.6cm, height=5.6cm,
    xmin=-1.25, xmax=1.25, ymin=-0.4, ymax=3.4,
    xtick={-1,1}, ytick={1.5708,3.1416},
    yticklabels={$\frac{\pi}{2}$,$\pi$}]
  \addplot[omThm, very thick, domain=-1:1, samples=200]
    {rad(acos(x))};
\end{axis}
\end{scope}
\begin{scope}[xshift=10.6cm]
\begin{axis}[omaxis, width=5.6cm, height=5.6cm,
    xmin=-4, xmax=4, ymin=-1.9, ymax=1.9,
    xtick={-2,2}, ytick={-1.5708,1.5708},
    yticklabels={$-\frac{\pi}{2}$,$\frac{\pi}{2}$}]
  \addplot[omMeth, very thick, domain=-4:4, samples=200]
    {rad(atan(x))};
  \addplot[dashed, gray, domain=-4:4] {1.5708};
  \addplot[dashed, gray, domain=-4:4] {-1.5708};
\end{axis}
\end{scope}
\end{tikzpicture}

{\small Dari kiri ke kanan: $\arcsin$, $\arccos$, $\arctan$. Masing-masing
mewarisi grafiknya dari fungsi langsung yang dibatasi lewat pencerminan
pada garis $y = x$.}
\end{omfigure}

\begin{proposition}[Turunannya]\label{prop:b1:functions:arcderiv}
Pada bagian dalam daerah asalnya:
\[
\arcsin' x = \frac{1}{\sqrt{1 - x^2}},
\qquad
\arccos' x = \frac{-1}{\sqrt{1 - x^2}},
\qquad
\arctan' x = \frac{1}{1 + x^2}.
\]
\end{proposition}

\begin{proof}
Dengan mendahului aturan fungsi invers yang dibuktikan pada
\cref{ch:b1:derivative}: jika $f$ bijeksi yang dapat diturunkan dan
$f' \neq 0$, maka $(f^{-1})'(x) = \frac{1}{f'(f^{-1}(x))}$. Untuk
$\arcsin$: dengan $y = \arcsin x$,
\[
\arcsin' x = \frac{1}{\cos y}
= \frac{1}{\sqrt{1 - \sin^2 y}} = \frac{1}{\sqrt{1 - x^2}},
\]
dengan $\cos y = +\sqrt{1 - \sin^2 y}$ karena $y \in
\intoo{-\frac\pi2}{\frac\pi2}$ memaksa $\cos y > 0$. Demikian pula
$\arccos' x = \frac{1}{-\sin y}$ dengan $\sin y > 0$ pada
$\intoo{0}{\pi}$, dan $\arctan' x = \frac{1}{1 + \tan^2 y} =
\frac{1}{1 + x^2}$ dengan memakai $\tan' = 1 + \tan^2$.
\end{proof}

\begin{example}[Mengenali konstanta yang tersembunyi]\label{ex:b1:functions:hiddenconst}
Pelajari $g(x) = \arctan\dfrac{1 - x}{1 + x}$ pada $\intoo{-1}
{+\infty}$. Aturan rantai dan sedikit perhitungan:
\[
g'(x) = \frac{1}{1 + \bigl(\frac{1-x}{1+x}\bigr)^2}\cdot
\frac{-(1+x) - (1-x)}{(1+x)^2}
= \frac{-2}{(1+x)^2 + (1-x)^2}
= \frac{-1}{1 + x^2} ,
\]
karena $(1+x)^2 + (1-x)^2 = 2 + 2x^2$. Jadi $g' = -\arctan'$: fungsi
$g + \arctan$ berturunan nol pada \emph{selang}
$\intoo{-1}{+\infty}$, sehingga konstan di sana; nilainya di $x =
0$ adalah $\arctan 1 + \arctan 0 = \frac\pi4$. Kesimpulannya:
\[
\arctan\frac{1 - x}{1 + x} = \frac\pi4 - \arctan x
\qquad (x > -1) .
\]
Pada $\intoo{-\infty}{-1}$ perhitungan turunan yang sama berlaku tetapi
konstantanya berbeda ($-\frac{3\pi}4$: hitung limitnya ketika
$x \to -\infty$). Inti gagasannya: ``turunan nol berarti
konstan'' adalah \omterm{def:b1:logic:statement}{pernyataan} selang demi selang --- persis
kehalusan yang dimanfaatkan pada \cref{exo:b1:functions:6} dan
\cref{exo:b1:functions:10}.
\end{example}

\begin{proposition}[Kesamaan baku]\label{prop:b1:functions:arcidentities}
\begin{enumerate}
  \item Untuk $x \in \intcc{-1}{1}$: $\arcsin x + \arccos x =
        \dfrac{\pi}{2}$.
  \item Untuk $x > 0$: $\arctan x + \arctan\dfrac 1x = \dfrac{\pi}{2}$
        (dan $-\frac\pi2$ untuk $x < 0$).
  \item $\sin(\arccos x) = \cos(\arcsin x) = \sqrt{1 - x^2}$;
        $\;\tan(\arcsin x) = \dfrac{x}{\sqrt{1 - x^2}}$ untuk
        $\abs x < 1$.
\end{enumerate}
\end{proposition}

\begin{proof}
(1) Turunan $x \mapsto \arcsin x + \arccos x$ bernilai nol pada
$\intoo{-1}{1}$ (\cref{prop:b1:functions:arcderiv}), jadi fungsinya
konstan di sana dan sama dengan nilainya $\frac\pi2$ di $0$; nilai di titik
ujung $\pm 1$ diperiksa langsung ($\frac\pi2 + 0$ dan $-\frac\pi2 +
\pi$).

(2) Metode yang sama pada $\intoo{0}{+\infty}$: turunannya adalah
$\frac{1}{1+x^2} + \frac{-1/x^2}{1 + 1/x^2} = \frac{1}{1+x^2} -
\frac{1}{x^2 + 1} = 0$, dan di $x = 1$ jumlahnya $2\arctan 1 =
\frac\pi2$. Untuk $x < 0$, pakai keganjilan $\arctan$.

(3) Dengan $y = \arccos x \in \intcc{0}{\pi}$: $\sin y \geq 0$, jadi
$\sin y = \sqrt{1 - \cos^2 y} = \sqrt{1 - x^2}$; demikian pula untuk
$\cos(\arcsin x)$, dan rumus tangennya adalah hasil baginya.
\end{proof}

\begin{remark}
$\arcsin(\sin\theta) = \theta$ berlaku \emph{hanya} untuk $\theta \in
\intcc{-\frac\pi2}{\frac\pi2}$: misalnya
$\arcsin(\sin\pi) = 0 \neq \pi$. Komposisi ke arah
sebaliknya, $\sin(\arcsin x) = x$, sah pada seluruh
$\intcc{-1}{1}$.
\end{remark}

\begin{example}[Melipat kembali ke selang utamanya]
\label{ex:b1:functions:wrap}
Hitunglah
\[
\arctan\Bigl(\tan\frac{3\pi}4\Bigr)
\qquad\text{dan}\qquad
\arctan\Bigl(\tan\frac{17\pi}5\Bigr).
\]
Resepnya: gantilah sudut
itu dengan satu-satunya sudut di $\intoo{-\frac\pi2}{\frac\pi2}$ yang
tangennya sama, yakni kurangkan kelipatan $\pi$ yang tepat
(yaitu periode $\tan$). Pertama: $\frac{3\pi}4 - \pi =
-\frac\pi4$, jadi jawabannya $-\frac\pi4$. Kedua: $\frac{17\pi}
5 - 3\pi = \frac{2\pi}5 \in \intoo{-\frac\pi2}{\frac\pi2}$, jadi
jawabannya $\frac{2\pi}5$. Perhitungan itu sesungguhnya pembagian
Euclid sudut itu oleh $\pi$ yang menyamar --- dan resep analognya
untuk $\arcsin$ (cerminkan ke dalam
$\intcc{-\frac\pi2}{\frac\pi2}$, dengan periode $2\pi$) dan $\arccos$
(cerminkan ke dalam $\intcc0\pi$) menggerakkan \cref{exo:b1:functions:1} dan
jawaban sepotong-sepotong pada \cref{exo:b1:functions:10}.
\end{example}

\begin{example}[Menjumlahkan arkus tangen dengan aman]\label{ex:b1:functions:arctansum}
Mari kita buktikan
\[
\arctan\frac12 + \arctan\frac15 + \arctan\frac18 = \frac\pi4 .
\]
Dua bahannya: rumus penjumlahan tangen, dan --- langkah yang
dilupakan pemula --- \emph{pelokalan} jumlahnya. Pertama,
$\tan(u + v) = \frac{\tan u + \tan v}{1 - \tan u\tan v}$ dengan
$u = \arctan\frac12$, $v = \arctan\frac15$ memberikan
\[
\tan(u + v) = \frac{\frac12 + \frac15}{1 - \frac1{10}}
= \frac{7/10}{9/10} = \frac79,
\qquad\text{lalu}\qquad
\tan\Bigl(u + v + \arctan\frac18\Bigr)
= \frac{\frac79 + \frac18}{1 - \frac7{72}}
= \frac{65/72}{65/72} = 1 .
\]
Kedua, masing-masing dari ketiga sudut itu terletak di $\intoo0{\frac\pi4}$
(karena argumennya kurang dari $1$), jadi jumlahnya terletak di
$\intoo0{\frac{3\pi}4}$; satu-satunya sudut di sana yang bertangen $1$ adalah
$\frac\pi4$. Tanpa pelokalan itu kesimpulannya hanya akan berbunyi
``$\frac\pi4$ sampai kelipatan $\pi$'' --- yaitu setengah bukti. Disiplin
dua langkah yang sama menggerakkan \cref{exo:b1:functions:8} dan
\cref{exo:b1:functions:12}.
\end{example}

\begin{example}[Sebuah persamaan arkus tangen]\label{ex:b1:functions:arctaneq}
Selesaikan $\arctan x + \arctan 2x = \dfrac\pi4$. Lokalkan lebih dulu:
ruas kirinya bertanda sama dengan $x$ (kedua sukunya begitu), jadi setiap penyelesaian
memenuhi $x > 0$, dan lalu jumlahnya terletak di $\intoo0\pi$. Ambil tangennya
(yang tak \omterm{def:b1:logic:inj}{injektif} pada selang mana pun yang panjangnya $\pi$, tetapi digabung dengan
pelokalan tadi ini sudah cukup): rumus penjumlahannya memberikan
\[
\tan\bigl(\arctan x + \arctan 2x\bigr)
= \frac{3x}{1 - 2x^2} = 1
\iff 2x^2 + 3x - 1 = 0
\iff x = \frac{-3 \pm \sqrt{17}}4 .
\]
Akar yang negatif dibuang oleh pelokalannya; untuk calon yang
positif $x_0 = \frac{\sqrt{17} - 3}4 \approx 0.28$,
jumlahnya terletak di $\intoo0\pi$ dengan tangen $1$, dan satu-satunya
sudut semacam itu adalah $\frac\pi4$ (pada $\intoo{\frac\pi2}\pi$ tangennya
negatif): jadi $x_0$ satu-satunya penyelesaian. Perhatikan bentuk
argumennya: mengambil tangen dapat melahirkan penyelesaian, tetapi tak pernah menghilangkannya,
jadi selesaikan dulu persamaan polinomialnya \emph{lalu} saring lewat
pelokalan --- disiplin periksa-maju yang sama seperti pada pengkuadratan
sebuah persamaan.
\end{example}

\section{Fungsi hiperbolik}

\begin{definition}[$\cosh$, $\sinh$, $\tanh$]\label{def:b1:functions:hyperbolic}
Untuk $x \in \R$:
\[
\cosh x = \frac{\eu^x + \eu^{-x}}{2},
\qquad
\sinh x = \frac{\eu^x - \eu^{-x}}{2},
\qquad
\tanh x = \frac{\sinh x}{\cosh x}
\]
(\emph{kosinus, sinus, tangen hiperbolik}\index{fungsi
hiperbolik}). Di sini $\cosh$ genap, sedangkan $\sinh$ dan $\tanh$ ganjil.
\end{definition}

\begin{proposition}[Sifat dasar]\label{prop:b1:functions:hyprules}
Untuk setiap $x, y \in \R$:
\begin{enumerate}
  \item $\cosh^2 x - \sinh^2 x = 1$;
  \item $\cosh' = \sinh$, $\;\sinh' = \cosh$, $\;\tanh' = 1 - \tanh^2
        = \dfrac{1}{\cosh^2}$;
  \item $\cosh(x+y) = \cosh x\cosh y + \sinh x \sinh y$ dan
        $\sinh(x+y) = \sinh x \cosh y + \cosh x \sinh y$;
  \item $\sinh$ adalah bijeksi naik tegas dari $\R$ pada $\R$;
        $\cosh$ yang dibatasi pada $\R_+$ adalah bijeksi naik
        tegas pada $\intco{1}{+\infty}$; dan $\tanh$ adalah bijeksi naik
        tegas dari $\R$ pada $\intoo{-1}{1}$.
\end{enumerate}
\end{proposition}

\begin{proof}
(1) $\bigl(\frac{\eu^x + \eu^{-x}}{2}\bigr)^2 - \bigl(\frac{\eu^x -
\eu^{-x}}{2}\bigr)^2 = \frac{(\eu^{2x} + 2 + \eu^{-2x}) - (\eu^{2x} -
2 + \eu^{-2x})}{4} = 1$.

(2) Turunkan rumus pendefinisinya; untuk $\tanh$, aturan hasil bagi
memberikan $\frac{\cosh^2 - \sinh^2}{\cosh^2}$, yang sekaligus
$1 - \tanh^2$ dan $\frac{1}{\cosh^2}$ menurut (1).

(3) Jabarkan ruas kanannya memakai definisinya. Secara rinci
untuk rumus yang pertama:
\[
\cosh x\cosh y + \sinh x\sinh y
= \frac{(\eu^x + \eu^{-x})(\eu^y + \eu^{-y})
+ (\eu^x - \eu^{-x})(\eu^y - \eu^{-y})}{4} ;
\]
dari kedelapan hasil kalinya, keempat yang ``campuran''
($\eu^{x-y}$, $\eu^{y-x}$) saling hapus berpasangan, sedangkan
$\eu^{x+y}$ dan $\eu^{-(x+y)}$ masing-masing muncul dua kali:
totalnya $\frac{2\eu^{x+y} + 2\eu^{-(x+y)}}{4} = \cosh(x+y)$.
Rumus sinusnya identik, hanya saja suku campurannya yang bertahan
sedangkan yang lain hapus.

(4) $\sinh' = \cosh \geq 1 > 0$, jadi $\sinh$ naik tegas,
dengan limit $\pm\infty$ (karena suku dominannya $\pm\frac{\eu^{\abs
x}}{2}$); \omterm{def:b1:logic:statement}{pernyataan} bijeksinya lalu menyusul dari teorema nilai
antara (yang dipakai pada taraf sekolah menengah; perlakuan sistematisnya pada
\cref{ch:b1:continuity}). Pada $\R_+$, $\cosh' = \sinh \geq 0$
dan hanya lenyap di $0$: jadi naik tegas dari $\cosh 0 = 1$ ke
$+\infty$. Dan $\tanh' > 0$ dengan limit $\pm 1$ di $\pm\infty$: secara
rinci,
\[
\tanh x = \frac{\eu^x - \eu^{-x}}{\eu^x + \eu^{-x}}
= \frac{1 - \eu^{-2x}}{1 + \eu^{-2x}}
\xrightarrow[x \to +\infty]{} 1 ,
\]
setelah membagi pembilang dan penyebutnya dengan $\eu^x$, sedangkan keganjilannya
memberikan limit $-1$ di $-\infty$; jadi fungsi $\tanh$ yang naik
tegas itu memetakan $\R$ pada $\intoo{-1}1$.
\end{proof}

\begin{omfigure}
\begin{tikzpicture}
\begin{axis}[omaxis, width=6.4cm, height=6cm,
    xmin=-2.6, xmax=2.6, ymin=-3.6, ymax=3.9,
    xtick={-2,-1,1,2}, ytick={-3,-2,-1,1,2,3}]
  \addplot[omDef, very thick, domain=-2.05:2.05, samples=120]
    {cosh(x)} node[pos=0.1, right, font=\small] {$\cosh$};
  \addplot[omThm, very thick, domain=-2.05:2.05, samples=120]
    {sinh(x)} node[pos=0.93, above left, font=\small] {$\sinh$};
  \addplot[dashed, gray, domain=0:2.5] {exp(x)/2};
\end{axis}
\begin{scope}[xshift=7.4cm]
\begin{axis}[omaxis, width=6.4cm, height=6cm,
    xmin=-3.2, xmax=3.2, ymin=-1.5, ymax=1.5,
    xtick={-2,-1,1,2}, ytick={-1,1}]
  \addplot[omMeth, very thick, domain=-3:3, samples=150]
    {tanh(x)} node[pos=0.9, above left, font=\small] {$\tanh$};
  \addplot[dashed, gray, domain=-3:3] {1};
  \addplot[dashed, gray, domain=-3:3] {-1};
\end{axis}
\end{scope}
\end{tikzpicture}

{\small Kiri: $\cosh$ dan $\sinh$, yang secara asimtotik merekat pada
$\frac{\eu^x}{2}$ (garis putus-putus). Kanan: $\tanh$, yang naik dari $-1$ ke
$1$.}
\end{omfigure}

\begin{remark}[Mengapa ``hiperbolik''?]
Titik $(\cosh t, \sinh t)$ menyusuri cabang $x > 0$ pada
hiperbola $x^2 - y^2 = 1$ (menurut
\cref{prop:b1:functions:hyprules}~(1)), persis seperti $(\cos t, \sin
t)$ menyusuri lingkaran $x^2 + y^2 = 1$. Setiap kesamaan lingkaran punya
saudara hiperboliknya, dengan perubahan tanda yang diatur oleh
$\sinh^2 \leftrightarrow -\sin^2$.
\end{remark}

\begin{example}[Rumus penjumlahan untuk $\tanh$]\label{ex:b1:functions:tanhadd}
Dengan membagi kedua rumus penjumlahan pada
\cref{prop:b1:functions:hyprules}~(3) dengan $\cosh x\cosh y$:
\[
\tanh(x + y)
= \frac{\sinh x\cosh y + \cosh x\sinh y}
{\cosh x\cosh y + \sinh x\sinh y}
= \frac{\tanh x + \tanh y}{1 + \tanh x\,\tanh y} ,
\]
yakni saudara hiperbolik dari rumus penjumlahan tangen --- dengan
$+$ yang di trigonometri berupa $-$. Sebuah dividen: karena $\abs{\tanh}
< 1$, ruas kanannya adalah kaidah ``penjumlahan kecepatan'' yang
tak pernah keluar dari $\intoo{-1}1$: jika $u, v \in \intoo{-1}1$ maka
$\frac{u + v}{1 + uv} \in \intoo{-1}1$ juga (tulis $u = \tanh
a$, $v = \tanh b$, yang mungkin karena kebijektifannya, lalu baca rumusnya
terbalik). Memeriksa hal itu secara aljabar, tanpa \omterm{def:b1:functions:hyperbolic}{fungsi
hiperbolik}, adalah latihan yang agak menyakitkan; memparameterkannya dengan
$\tanh$ membuatnya satu baris --- strategi yang sama seperti yang
diberikan fungsi lingkaran bagi lingkaran satuan.
\end{example}

\begin{proposition}[Fungsi invers hiperbolik]\label{prop:b1:functions:invhyp}
Invers yang dijamin oleh \cref{prop:b1:functions:hyprules}~(4) mempunyai
bentuk tertutup:
\[
\operatorname{arsinh} x = \ln\bigl(x + \sqrt{x^2 + 1}\bigr)
\ (x \in \R),
\qquad
\operatorname{arcosh} x = \ln\bigl(x + \sqrt{x^2 - 1}\bigr)
\ (x \geq 1),
\]
\[
\operatorname{artanh} x = \frac 12 \ln\frac{1 + x}{1 - x}
\ (\abs x < 1),
\]
dengan turunan berturut-turut $\frac{1}{\sqrt{x^2+1}}$, $\frac{1}{\sqrt{x^2-1}}$
($x > 1$) dan $\frac{1}{1 - x^2}$.
\end{proposition}

\begin{proof}
Untuk $\operatorname{arsinh}$: selesaikan $x = \sinh y = \frac{\eu^y -
\eu^{-y}}{2}$. Dengan menulis $u = \eu^y > 0$: $u^2 - 2xu - 1 = 0$, jadi $u = x
+ \sqrt{x^2 + 1}$ (karena akar $x - \sqrt{x^2+1}$ negatif), dan $y =
\ln(x + \sqrt{x^2+1})$. Dua yang lain adalah perhitungan yang identik
($u^2 - 2xu + 1 = 0$ untuk $\operatorname{arcosh}$, dengan mempertahankan akar
yang $\geq 1$; dan penataan ulang dua baris untuk $\operatorname{artanh}$).
Turunannya: turunkan ungkapan logaritmanya, misalnya
\[
\bigl(\ln(x + \sqrt{x^2+1})\bigr)'
= \frac{1 + \frac{x}{\sqrt{x^2+1}}}{x + \sqrt{x^2+1}}
= \frac{1}{\sqrt{x^2 + 1}} . \qedhere
\]
\end{proof}

\begin{example}[Bentuk tertutup dalam praktik]\label{ex:b1:functions:arcoshtwo}
Penyelesaian $\cosh t = 2$ dengan $t \geq 0$ adalah, menurut bentuk
tertutupnya, $t = \operatorname{arcosh} 2 = \ln(2 + \sqrt3) \approx
1.317$; penyelesaian yang lain adalah $-t$, karena kegenapannya --- dan memang
$\ln(2 - \sqrt3) = \ln\frac1{2 + \sqrt3} = -\ln(2 + \sqrt3)$: kedua
akar $u = \eu^t$ persamaan kuadrat $u^2 - 4u + 1 = 0$ saling
berkebalikan, sebagaimana dituntut hasil kalinya $1$ (menurut Vieta). Perhitungan
kecil ini memperlihatkan pola umumnya: persamaan hiperbolik
berubah menjadi persamaan kuadrat dalam $\eu^t$, dan simetri $t \mapsto -t$
muncul sebagai simetri $u \mapsto 1/u$ persamaan kuadrat itu ---
layak diingat ketika menyelesaikan \cref{exo:b1:functions:7}.
\end{example}

\begin{method}[Memilih bentuk antiturunan yang tepat]\label{met:b1:functions:primitive}
Ketiga pola turunan yang perlu dihafalkan untuk pengintegralan
(\cref{ch:b1:integration}):
\[
\int \frac{\dd x}{1 + x^2} = \arctan x + C,
\qquad
\int \frac{\dd x}{\sqrt{1 - x^2}} = \arcsin x + C,
\qquad
\int \frac{\dd x}{\sqrt{x^2 + 1}} = \operatorname{arsinh} x + C,
\]
dan untuk kuadrat yang umum, susutkan ke bentuk itu dengan melengkapkan kuadratnya
lalu menskalakannya.
\end{method}

\begin{remark}[Di mana bab ini dipakai]\label{rem:b1:functions:whereused}
Bab ini adalah kosakata kerja bagi seluruh analisis yang akan datang.
Skala pertumbuhan pada \cref{prop:b1:functions:powerrules} memutuskan
pertanyaan kekonvergenan di sepanjang
\cref{ch:b1:seq,ch:b1:series}; rumus turunan pada
\cref{prop:b1:functions:arcderiv} dan
\cref{prop:b1:functions:invhyp} adalah antiturunan yang paling sering
diperlukan pada \cref{ch:b1:integration}, lewat
\cref{met:b1:functions:primitive}; \omterm{def:b1:functions:hyperbolic}{fungsi hiperbolik}
memparameterkan penyelesaian persamaan $y'' = y$ pada
\cref{ch:b1:diffeq} persis seperti fungsi lingkaran memparameterkan
penyelesaian $y'' = -y$. Pencirian $\exp$ lewat $f' = f$,
$f(0) = 1$ (\cref{prop:b1:functions:expln}) adalah benih berdimensi satu
bagi teori persamaan diferensial linear, dan
kurva rantai $y = \cosh x$ kembali di antara kurva bidang pada
\cref{ch:b1:curves}.
\end{remark}

\begin{remark}[Selingan: program fungsi invers]\label{rem:b1:functions:interlude}
Bab ini menjalankan satu program sebanyak empat kali: batasi sebuah fungsi
sampai ia menjadi bijeksi yang monoton tegas, namai inversnya,
lalu angkut setiap rumus melewatinya. Apa yang \emph{dipakai} pada masing-masing
langkah --- bahwa fungsi kontinu yang monoton tegas pada sebuah
selang merupakan bijeksi pada sebuah selang, dan bahwa inversnya
kontinu, lalu dapat diturunkan di luar titik kritisnya --- dipinjam
atas kredit sekolah menengah. Utang itu dilunasi di dalam jilid
ini: \cref{ch:b1:continuity} membuktikan \omterm{def:b1:logic:statement}{pernyataan} bijeksinya
(teorema invers monoton, yang bersandar pada teorema nilai
antara), dan \cref{ch:b1:derivative} membuktikan aturan turunan
$(f^{-1})' = 1/(f' \circ f^{-1})$ yang diam-diam menghasilkan setiap
rumus pada \cref{prop:b1:functions:arcderiv} dan
\cref{prop:b1:functions:invhyp}. Membaca bab-bab itu dengan
$\arcsin$ dan $\operatorname{arsinh}$ di kepala --- sebagai contoh
terkerjakan yang untuknyalah teorinya dibangun --- adalah cara yang
dimaksudkan untuk mengatasi kemelingkaran yang tampak itu.
\end{remark}

\section{Latihan}

\begin{exercise}[$\star$]\label{exo:b1:functions:1}
Hitung tanpa kalkulator: $\arcsin\frac{\sqrt 3}{2}$;
$\;\arccos\bigl(-\frac 12\bigr)$; $\;\arctan(-1)$;
$\;\arcsin\bigl(\sin\frac{5\pi}{6}\bigr)$;
$\;\arccos\bigl(\cos\frac{7\pi}{4}\bigr)$.
\end{exercise}

\begin{exercise}[$\star$]\label{exo:b1:functions:2}
Berikan daerah definisi $f(x) = \arcsin(2x - 1)$ lalu hitung
$f'$ di tempat ia terdefinisi. Pertanyaan yang sama untuk $g(x) = \arctan\sqrt{x}$.
\end{exercise}

\begin{exercise}[$\star$]\label{exo:b1:functions:3}
Buktikan bahwa untuk setiap $x \in \R$: $\cosh x + \sinh x = \eu^x$ dan
$(\cosh x + \sinh x)^n = \cosh nx + \sinh nx$ untuk setiap $n \in \Z$
(yakni rumus de Moivre yang hiperbolik).
\end{exercise}

\begin{exercise}[$\star$]\label{exo:b1:functions:4}
Sederhanakan $\cos(2\arcsin x)$ dan $\sin(2\arctan x)$ menjadi ungkapan
aljabar dalam $x$.
\end{exercise}

\begin{exercise}[$\star$]\label{exo:b1:functions:5}
Urutkan, untuk $x$ yang besar, fungsi
$x^{100}$, $\;\eu^{\sqrt x}$, $\;(\ln x)^{1000}$, $\;\eu^{x/100}$,
$\;x^{\ln x}$, dari yang pertumbuhannya paling lambat ke paling cepat, dengan alasan
yang berdasar pada \cref{prop:b1:functions:powerrules}.
\end{exercise}

\begin{exercise}[$\star\star$]\label{exo:b1:functions:6}
Pelajari fungsi $f(x) = \arctan\dfrac{2x}{1 - x^2}$ pada daerah asalnya:
hitung $f'$, bandingkan dengan $(2\arctan x)'$, lalu nyatakan $f(x)$
dengan $\arctan x$ pada masing-masing dari ketiga selang daerah asalnya.
\end{exercise}

\begin{exercise}[$\star\star$]\label{exo:b1:functions:7}
Selesaikan di $\R$: $\;\cosh x = 2$; lalu $5\cosh x - 4 \sinh x = 3$
(nyatakan penyelesaiannya dengan logaritma). \emph{Petunjuk untuk yang kedua:
tuliskan semuanya dengan $u = \eu^x$.}
\end{exercise}

\begin{exercise}[$\star\star$]\label{exo:b1:functions:8}
Buktikan kesamaan $\arctan\frac{1}{2} + \arctan\frac{1}{3} =
\frac{\pi}{4}$, lalu rumus Machin\index{rumus Machin}
\[
4\arctan\frac 15 - \arctan\frac{1}{239} = \frac{\pi}{4}.
\]
\emph{Petunjuk: hitung tangen kedua ruasnya memakai rumus
penjumlahan, lalu kendalikan di selang mana masing-masing ruas terletak.}
\end{exercise}

\begin{exercise}[$\star\star$]\label{exo:b1:functions:9}
Buktikan bahwa untuk setiap $x \geq 0$: $\;x - \dfrac{x^3}{6} \leq \sin x \leq
x$ \emph{(pelajari turunan berturutan dari selisihnya)}, lalu
simpulkan bahwa $\lim_{x \to 0^+} \frac{x - \sin x}{x^3} = \frac 16$
masuk akal --- limitnya sendiri ditegakkan pada
\cref{ch:b1:taylor}.
\end{exercise}

\begin{exercise}[$\star\star\star$]\label{exo:b1:functions:10}
Untuk $x \in \intcc{-1}{1}$, tulis
$h(x) = \arcsin\bigl(2x\sqrt{1 - x^2}\,\bigr) - 2\arcsin x$. Orang mungkin
mengharapkan $h = 0$ dari kesamaan $\sin 2y = 2\sin y\cos y$ --- tetapi
$h$ \emph{tidak} identik nol. Hitung $h'$ pada selang terbuka
yang di situ ia ada, hitung nilai $h$ pada titik yang dipilih baik, lalu berikan
uraian lengkap $h$ sebagai fungsi konstan sepotong-sepotong pada $\intcc{-1}{1}$.
\end{exercise}

\begin{exercise}[$\star\star\star$]\label{exo:b1:functions:11}
(Fungsi Gudermann) Misalkan $g(x) = \arctan(\sinh x)$ untuk $x \in \R$. Buktikan
bahwa $g$ adalah bijeksi ganjil yang naik tegas dari $\R$ pada
$\intoo{-\frac\pi2}{\frac\pi2}$, bahwa $g'(x) = \frac{1}{\cosh x}$,
dan bahwa $\tan g(x) = \sinh x$, $\;\sin g(x) = \tanh x$,
$\;\cos g(x) = \frac{1}{\cosh x}$: jadi fungsi $g$ merangkai trigonometri
lingkaran dan hiperbolik tanpa bilangan kompleks.
\end{exercise}

\begin{exercise}[$\star\star$]\label{exo:b1:functions:12}
Buktikan bahwa
\[
\arctan 1 + \arctan 2 + \arctan 3 = \pi .
\]
\emph{Petunjuk: hitung $\arctan 2 + \arctan 3$ lebih dulu, dengan melokalkan
jumlahnya seperti pada \cref{ex:b1:functions:arctansum}; awas, penyebut
rumus penjumlahannya bernilai negatif di sini.}
\end{exercise}

\section{Soal: Persamaan $x^y = y^x$}

\begin{problem}\label{pb:b1:functions:1}
Pasangan bilangan positif mana yang memenuhi $x^y = y^x$? Semua orang tahu
satu contoh yang tampak kebetulan, $2^4 = 4^2 = 16$; soal ini menunjukkan
bahwa tak ada yang kebetulan di situ. Seluruh persamaannya dikendalikan
oleh variasi satu fungsi tunggal
\[
f(t) = \frac{\ln t}{t} \qquad (t > 0),
\]
yang penelaahannya menghasilkan: \omterm{def:b1:logic:sets}{himpunan} penyelesaian lengkapnya (sebuah diagonal ditambah satu
cabang melengkung yang melewati $(\eu, \eu)$), sebuah parameterisasi rasional
cabang itu, kenyataan bahwa $(2, 4)$ adalah \emph{satu-satunya} titik bulatnya
dan penggolongan semua titik rasionalnya --- ditambah, sebagai
dividen, perbandingan $\eu^\pi$ dengan $\pi^\eu$ dan
kekonvergenan monoton $\bigl(1 + \frac1n\bigr)^n$ ke $\eu$.
\index{fungsi pangkat}

\medskip
\noindent\textbf{Bagian I --- Fungsi $f(t) = \ln t / t$.}
\begin{enumerate}
  \item Berikan alasan bahwa $f$ dapat diturunkan pada $\intoo0{+\infty}$,
        hitung $f'$, lalu susun tabel variasinya: $f$
        naik tegas pada $\intoc0\eu$, turun tegas pada
        $\intco\eu{+\infty}$, dengan maksimum $f(\eu) = \frac1\eu$.
  \item Tentukan limit $f$ di $0^+$ dan di $+\infty$
        (\cref{prop:b1:functions:powerrules}), tanda $f$
        (negatif pada $\intoo01$, nol di $1$, positif setelah itu),
        lalu sketsakan grafiknya.
  \item Periksa lewat perhitungan langsung bahwa $f(2) = f(4)$. (Simpanlah
        kesamaan ini di kepala: seluruh soal ini tumbuh darinya.)
  \item Misalkan $c \in \R$. Bahaslah, menurut nilai $c$, banyaknya
        penyelesaian $f(t) = c$: tepat satu untuk $c \leq
        0$; tepat dua (satu di $\intoo1\eu$, satu di
        $\intoo\eu{+\infty}$) untuk $0 < c < \frac1\eu$; tepat satu
        untuk $c = \frac1\eu$; dan tak ada untuk $c > \frac1\eu$.
  \item Simpulkan bahwa $\eu^x > x^\eu$ untuk setiap $x > 0$ dengan $x \neq
        \eu$, lalu khususnya tentukan mana yang lebih besar di antara $\eu^\pi$ dan
        $\pi^\eu$.
\end{enumerate}

\medskip
\noindent\textbf{Bagian II --- Persamaan itu dan kurvanya.} Pada bagian
ini $x, y > 0$.
\begin{enumerate}[resume]
  \item Tunjukkan bahwa $x^y = y^x \iff f(x) = f(y)$.
  \item Simpulkan struktur \omterm{def:b1:logic:sets}{himpunan} penyelesaiannya: semua pasangan
        diagonal $(x, x)$; dan pasangan yang \emph{tak sepele} ($x \neq
        y$), yang memenuhi: kedua koordinatnya $> 1$, dan satu di antaranya
        terletak di $\intoo1\eu$ sedangkan yang lain terletak di
        $\intoo\eu{+\infty}$.
  \item Parameterkan pasangan yang tak sepele itu: dengan menulis $y = tx$ dengan $t
        > 0$, $t \neq 1$, tunjukkan bahwa $x^y = y^x$ memaksa
        \[
        x(t) = t^{\frac1{t-1}},
        \qquad
        y(t) = t\,x(t) = t^{\frac{t}{t-1}},
        \]
        dan bahwa sebaliknya setiap pasangan semacam itu merupakan penyelesaian.
  \item Periksa bahwa $t = 2$ memberikan $(2, 4)$, lalu buktikan simetri
        $x(1/t) = y(t)$, $y(1/t) = x(t)$: membalikkan parameternya
        menukarkan kedua koordinatnya.
  \item Tentukan limit $x(t)$ dan $y(t)$ ketika $t \to 1$
        (keduanya menuju $\eu$), ketika $t \to +\infty$ ($x \to 1$, $y
        \to +\infty$) dan ketika $t \to 0^+$ ($x \to +\infty$, $y \to
        1$). Gambarkan cabang yang dihasilkannya: sebuah kurva yang asimtotik ke
        garis $x = 1$ dan $y = 1$, yang memotong diagonalnya di
        $(\eu, \eu)$.
  \item Tunjukkan bahwa $t \mapsto x(t)$ turun tegas pada
        $\intoo1{+\infty}$. \emph{(Pelajari $h(t) = 1 - \frac1t - \ln
        t$, yang tandanya mengendalikan turunan $\ln x(t) =
        \frac{\ln t}{t - 1}$.)}
  \item Simpulkan bahwa penyelesaian yang tak sepele mendefinisikan bijeksi
        turun tegas $\varphi \colon \intoo1\eu \to
        \intoo\eu{+\infty}$ dengan $x^{\varphi(x)} = \varphi(x)^x$,
        dan bahwa $\varphi$ merupakan involusi pada cabang itu:
        $\varphi(\varphi(x)) = x$ di mana pun kedua ruasnya terdefinisi.
\end{enumerate}

\medskip
\noindent\textbf{Bagian III --- Titik bulat dan titik rasional.}
\begin{enumerate}[resume]
  \item Buktikan bahwa $(2, 4)$ dan $(4, 2)$ adalah satu-satunya penyelesaian
        \emph{bulat} yang tak sepele bagi $x^y = y^x$.
  \item Untuk $n \in \N^*$, terapkan parameterisasinya dengan $t = 1 +
        \frac1n$ lalu tunjukkan bahwa
        \[
        x_n = \Bigl(1 + \frac1n\Bigr)^{n},
        \qquad
        y_n = \Bigl(1 + \frac1n\Bigr)^{n+1}
        \]
        merupakan penyelesaian \emph{rasional} yang tak sepele untuk setiap $n$,
        dengan $(x_1, y_1) = (2, 4)$.
  \item Sebaliknya, misalkan $(x, y)$ penyelesaian rasional yang tak sepele
        dengan $y > x$, lalu tulis $t = y/x = 1 + \frac rs$ dalam bentuk paling
        sederhana ($r, s \in \N^*$). Dengan memakai $x = t^{s/r}$ dan menerima
        ketunggalan pemfaktoran prima (yang dikenal dari
        sekolah; dibuktikan pada \cref{ch:b1:arith}), tunjukkan bahwa $x$
        yang rasional memaksa $s$ maupun $s + r$ menjadi pangkat ke-$r$
        sebuah bilangan bulat.
  \item Tunjukkan dengan teorema binomial bahwa $a^r - b^r = r$ tak mempunyai
        penyelesaian bulat $a > b \geq 1$ ketika $r \geq 2$, lalu
        simpulkan: penyelesaian rasional $x^y = y^x$ persis
        pasangan $(x_n, y_n)$ pada pertanyaan 14 (beserta
        tukarannya).
  \item Periksa kasus $n = 2$ secara numerik: hitung $f(9/4)$ dan
        $f(27/8)$ sampai empat angka desimal lalu periksa bahwa keduanya cocok.
\end{enumerate}

\medskip
\noindent\textbf{Bagian IV --- Dividennya.}
\begin{enumerate}[resume]
  \item Pakai pertanyaan 7 dan 11 untuk membuktikan, tanpa perhitungan
        tambahan, bahwa barisan $x_n = \bigl(1 +
        \frac1n\bigr)^n$ naik tegas dengan $x_n < \eu$,
        bahwa $y_n = \bigl(1 + \frac1n\bigr)^{n+1}$ turun
        tegas dengan $y_n > \eu$, dan bahwa keduanya konvergen ke
        $\eu$. (Parameternya $t_n = 1 + \frac1n$ turun ke
        $1$.)
  \item Tegakkan kaidah perbandingan umum untuk $1 < a < b$:
        jika $\eu \leq a$ maka $a^b > b^a$; jika $b \leq \eu$ maka
        $a^b < b^a$; lalu tunjukkan lewat dua contoh $(2, 3)$ dan
        $(2, 5)$ bahwa pada kasus campuran $a < \eu < b$ kedua
        hasil itu sungguh-sungguh terjadi.
  \item Andaikan involusi $\varphi$ pada pertanyaan 12 dapat
        diturunkan di $\eu$ (memang begitu). Turunkan
        kesamaan $\varphi(\varphi(x)) = x$ di $x = \eu$ lalu simpulkan
        $\varphi'(\eu) = -1$: jadi cabang itu memotong diagonalnya
        secara tegak lurus.
  \item Carilah satu-satunya pasangan penyelesaian dengan $y = 3x$, dalam bentuk
        tertutup, lalu periksa secara numerik sampai empat angka desimal lewat $f$.
  \item Di antara bilangan $\sqrt2$, $\sqrt[3]3$, $\sqrt[4]4$,
        $\sqrt[5]5$, tentukan yang terbesar lalu kenali dua di antaranya
        yang sama. \emph{(Bandingkan $n^{1/n} =
        \eu^{f(n)}$.)}
\end{enumerate}

\medskip
\noindent\textbf{Bagian V --- Sintesis.}
\begin{enumerate}[resume]
  \item Gambarkan \omterm{def:b1:logic:sets}{himpunan} penyelesaian lengkap $x^y = y^x$ pada
        seperempat bidang $x, y > 0$ --- diagonal ditambah cabangnya, titik
        potongnya $(\eu, \eu)$, asimtotnya, titik bulat
        $(2, 4)$, titik rasional yang menumpuk di
        $(\eu, \eu)$ --- dalam bentuk yang dapat Anda sketsakan dari ingatan.
  \item Di mana persisnya soal ini memakai: (i) \omterm{prop:b1:functions:powerrules}{perbandingan
        pertumbuhan} pada \cref{prop:b1:functions:powerrules}; (ii)
        teorema nilai antara (lewat \omterm{def:b1:logic:statement}{pernyataan}
        bijeksinya); (iii) ketunggalan pemfaktoran prima yang
        diterima tanpa bukti? Satu kalimat untuk masing-masing.
  \item Moralnya, dalam satu paragraf pendek: satu tabel variasi
        menyelesaikan sebuah persamaan dengan dua variabel, menggolongkan titik
        rasionalnya, dan membuktikan kekonvergenan monoton
        $\bigl(1 + \frac1n\bigr)^n$ --- berilah komentar atas penghematan ini,
        lalu sebutkan tempat masing-masing benangnya diindustrialisasi di bagian
        akhir jilid ini (\cref{ch:b1:seq} untuk barisannya,
        \cref{ch:b1:derivative} untuk tabel variasinya,
        \cref{ch:b1:taylor} untuk ketelitian yang tak dimiliki tabel itu).

\end{enumerate}
\end{problem}
